% Figure 1 (fig:overview) — TikZ source. This file is \input by
% sections/03_methodology.tex inside the figure* environment; iterate HERE.
%
% v2 (redesign on Robert's notes, 2026-07-12):
%   1. QUANTIZATION IS VISIBLE: every pickup/taxi mark is snapped to a cell
%      center and sized to fill most of its cell — the grid is the data model.
%   2. CITY UNDER THE GRID: a stylized street network (gray arterials + minor
%      roads, cGAIL-figure style) drawn natively in TikZ under each grid, plus
%      a district boundary (dash, follows cell edges) splitting an advantaged
%      district (white) from a disadvantaged one (light gray tint) — districts
%      carry the demographics semantics.
%   3. Trim edits are vermillion arrows with a legend; Panel 3 is a 2x detail
%      view of the under-served district (zoom box marked in Panel 2) showing
%      BOTH the original tail and the rerouted tail, with a legend; the
%      value-of-presence glow is now QUANTIZED (per-cell accent shading).
%
% Color scheme (supersedes the one-accent note in PAPER/figures/figure-1.md):
%   TWO hues, colorblind-safe pair — blue #0F62FE = lift/supply-side,
%   vermillion #D55E00 = trim/demand-side edits. Both redundant with shape
%   (dash pattern / arrowheads), so meaning survives grayscale. Roads,
%   districts, and marks stay in the gray ramp.
%
% Geometry: Panels 1-2 = the full 10x8-cell city (1 cell = 1 unit = 4.9mm).
% Panel 3 = the city sub-window x in [5,10], y in [0,4] at 2x zoom (1 city
% cell = 2 panel units), so cells visibly enlarge — honest zoom semantics.
% No \resizebox: fonts are true manuscript size.
\definecolor{figink}{HTML}{1A1A1A}
\definecolor{figgray}{HTML}{6F6F6F}
\definecolor{figfaint}{HTML}{C9C9C9}
\definecolor{figaccent}{HTML}{0F62FE}
\definecolor{figtrim}{HTML}{D55E00}
\definecolor{figroad}{HTML}{BFBFBF}
\definecolor{figminor}{HTML}{DCDCDC}
\begin{tikzpicture}[x=4.9mm, y=4.9mm, line cap=round, line join=round,
  % cell-scale marks (panels 1-2): \pic[color=<c>] at (cell center) {pick};
  pick/.pic={\draw[line width=0.9pt] (-0.26,-0.26) -- (0.26,0.26)
             (-0.26,0.26) -- (0.26,-0.26);},
  taxi/.pic={\draw[figink, fill=white, line width=0.7pt]
             (-0.26,-0.26) rectangle (0.26,0.26);},
  % 2x marks for the Panel-3 detail view
  pickz/.pic={\draw[line width=1.1pt] (-0.5,-0.5) -- (0.5,0.5)
              (-0.5,0.5) -- (0.5,-0.5);},
  taxiz/.pic={\draw[figink, fill=white, line width=0.9pt]
              (-0.5,-0.5) rectangle (0.5,0.5);},
  trimedit/.style={-{Stealth[length=3.6pt]}, figtrim, line width=1.1pt},
  liftdash/.style={figaccent, dashed, line width=1.1pt},
  boundary/.style={figink!65, line width=1pt, dash pattern=on 4pt off 2pt},
  arterial/.style={figroad, line width=1.5pt, rounded corners=3pt},
  minorroad/.style={figminor, line width=0.5pt},
  panelframe/.style={figgray, line width=0.6pt},
  citygrid/.style={figfaint, line width=0.4pt}]

  % =========================================================================
  % Panel 1: the service gap  (full city)
  % =========================================================================
  \begin{scope}
    % district tint (disadvantaged, east of the boundary), then roads, grid
    \fill[black!7] (4,8) -- (4,6) -- (5,6) -- (5,4) -- (6,4) -- (6,2)
      -- (7,2) -- (7,0) -- (10,0) -- (10,8) -- cycle;
    \draw[arterial] (0,6.3) -- (2,6.1) -- (4.5,6.4) -- (7,5.8) -- (10,6.0);
    \draw[arterial] (2.8,0) -- (3.1,2.5) -- (2.7,5) -- (3.0,8);
    \draw[arterial] (0,2.2) -- (3,2.6) -- (6,1.8) -- (10,2.4);
    \draw[arterial] (7.6,8) -- (7.2,5.5) -- (7.8,3) -- (7.5,0);
    \draw[minorroad] (0,4.1) -- (4.5,4.0)  (5.5,4.9) -- (10,5.1)
      (1.2,8) -- (1.4,4.2)  (8.9,8) -- (9.1,4.5)  (5.6,0) -- (5.8,3.2);
    \draw[citygrid, step=1] (0,0) grid (10,8);
    \draw[boundary] (4,8) -- (4,6) -- (5,6) -- (5,4) -- (6,4) -- (6,2)
      -- (7,2) -- (7,0);
    \draw[panelframe] (0,0) rectangle (10,8);
    \node[anchor=south west, font=\footnotesize, text=figink] at (0,8.2)
      {1\,$\cdot$\,the service gap};
    % advantaged district: dense presence (cell-snapped), few pickups
    \foreach \x/\y in {1.5/6.5, 2.5/6.5, 3.5/6.5, 1.5/5.5, 3.5/5.5,
                       1.5/4.5, 2.5/4.5, 3.5/4.5, 2.5/7.5, 0.5/5.5}
      \pic at (\x,\y) {taxi};
    \foreach \x/\y in {2.5/5.5, 0.5/6.5, 0.5/4.5, 3.5/7.5}
      \pic[color=figgray] at (\x,\y) {pick};
    % disadvantaged district: dense unmet pickups, one taxi
    \foreach \x/\y in {8.5/1.5, 7.5/1.5, 8.5/2.5, 9.5/1.5, 7.5/0.5,
                       8.5/3.5, 9.5/2.5, 9.5/3.5}
      \pic[color=figgray] at (\x,\y) {pick};
    \pic at (7.5,2.5) {taxi};
    % district + service labels
    \node[font=\tiny, text=figink, align=center] at (2.1,3.4)
      {advantaged district\\(over-served)};
    \node[font=\tiny, text=figink, align=center] at (8.3,4.9)
      {disadvantaged district\\(under-served)};
    % legend (lower left, advantaged side is clear there)
    \pic at (0.5,1.05) {taxi};
    \node[anchor=west, font=\tiny, text=figink] at (0.85,1.05) {taxi presence};
    \pic[color=figgray] at (0.5,0.4) {pick};
    \node[anchor=west, font=\tiny, text=figink] at (0.85,0.4) {pickup};
  \end{scope}

  % =========================================================================
  % Panel 2: trim  (same city; two pickups relocated inside the district)
  % =========================================================================
  \begin{scope}[shift={(11.5,0)}]
    \fill[black!7] (4,8) -- (4,6) -- (5,6) -- (5,4) -- (6,4) -- (6,2)
      -- (7,2) -- (7,0) -- (10,0) -- (10,8) -- cycle;
    \draw[arterial] (0,6.3) -- (2,6.1) -- (4.5,6.4) -- (7,5.8) -- (10,6.0);
    \draw[arterial] (2.8,0) -- (3.1,2.5) -- (2.7,5) -- (3.0,8);
    \draw[arterial] (0,2.2) -- (3,2.6) -- (6,1.8) -- (10,2.4);
    \draw[arterial] (7.6,8) -- (7.2,5.5) -- (7.8,3) -- (7.5,0);
    \draw[minorroad] (0,4.1) -- (4.5,4.0)  (5.5,4.9) -- (10,5.1)
      (1.2,8) -- (1.4,4.2)  (8.9,8) -- (9.1,4.5)  (5.6,0) -- (5.8,3.2);
    \draw[citygrid, step=1] (0,0) grid (10,8);
    \draw[boundary] (4,8) -- (4,6) -- (5,6) -- (5,4) -- (6,4) -- (6,2)
      -- (7,2) -- (7,0);
    \draw[panelframe] (0,0) rectangle (10,8);
    \node[anchor=south west, font=\footnotesize, text=figink] at (0,8.2)
      {2\,$\cdot$\,trim: relocate excess pickups};
    % same taxis; the two unedited pickups stay put
    \foreach \x/\y in {1.5/6.5, 2.5/6.5, 3.5/6.5, 1.5/5.5, 3.5/5.5,
                       1.5/4.5, 2.5/4.5, 3.5/4.5, 2.5/7.5, 0.5/5.5}
      \pic at (\x,\y) {taxi};
    \foreach \x/\y in {0.5/6.5, 3.5/7.5}
      \pic[color=figgray] at (\x,\y) {pick};
    % untouched disadvantaged district — identical to Panel 1
    \foreach \x/\y in {8.5/1.5, 7.5/1.5, 8.5/2.5, 9.5/1.5, 7.5/0.5,
                       8.5/3.5, 9.5/2.5, 9.5/3.5}
      \pic[color=figgray] at (\x,\y) {pick};
    \pic at (7.5,2.5) {taxi};
    % two trim edits: ghosted origin -> vermillion arrow -> relocated pickup,
    % both staying inside the advantaged district (west of the boundary)
    \pic[color=figfaint] at (2.5,5.5) {pick};
    \draw[trimedit] (2.5,5.5) -- (4.5,4.5);
    \pic[color=figgray] at (4.5,4.5) {pick};
    \pic[color=figfaint] at (0.5,4.5) {pick};
    \draw[trimedit] (0.5,4.5) -- (1.5,3.5);
    \pic[color=figgray] at (1.5,3.5) {pick};
    \node[font=\tiny, text=figgray, align=center] at (8.0,6.9)
      {levels down:\\under-served untouched};
    % zoom cue: Panel 3 shows this window in detail
    \draw[figink, line width=0.6pt, dash pattern=on 2pt off 1.6pt]
      (5,0) rectangle (10,4);
    \node[anchor=south east, font=\tiny, text=figink] at (10,4.1) {panel 3};
    % legend (lower left)
    \pic[color=figfaint] at (0.5,1.05) {pick};
    \node[anchor=west, font=\tiny, text=figink] at (0.85,1.05)
      {original position};
    \draw[trimedit] (0.25,0.4) -- (0.8,0.4);
    \node[anchor=west, font=\tiny, text=figink] at (0.85,0.4)
      {trim edit (pickup relocated)};
  \end{scope}

  % =========================================================================
  % Panel 3: lift  (2x detail of the under-served district: city [5,10]x[0,4];
  % panel coords = 2*(city - (5,0)), so one city cell = a 2x2 block here)
  % =========================================================================
  \begin{scope}[shift={(23,0)}]
    % district tint east of the boundary (continues Panels 1-2)
    \fill[black!7] (2,8) -- (2,4) -- (4,4) -- (4,0) -- (10,0) -- (10,8)
      -- cycle;
    % streets at 2x (portions of the same arterials)
    \draw[arterial, line width=2.6pt] (0,4.0) -- (4,4.6) -- (8,3.4) -- (10,4.4);
    \draw[arterial, line width=2.6pt] (5.4,8) -- (4.8,5) -- (5.6,2) -- (5.2,0);
    \draw[minorroad, line width=0.7pt] (1.2,0) -- (1.6,8);
    % QUANTIZED value-of-presence field (accent cells, opacity = value)
    \fill[figaccent, fill opacity=0.26] (6,2) rectangle (8,4);   % peak
    \fill[figaccent, fill opacity=0.15] (4,2) rectangle (6,4);
    \fill[figaccent, fill opacity=0.15] (8,2) rectangle (10,4);
    \fill[figaccent, fill opacity=0.15] (6,4) rectangle (8,6);
    \fill[figaccent, fill opacity=0.15] (6,0) rectangle (8,2);
    \fill[figaccent, fill opacity=0.07] (4,0) rectangle (6,2);
    \fill[figaccent, fill opacity=0.07] (8,0) rectangle (10,2);
    \fill[figaccent, fill opacity=0.07] (8,4) rectangle (10,6);
    \draw[citygrid, step=2] (0,0) grid (10,8);
    \draw[boundary] (2,8) -- (2,4) -- (4,4) -- (4,0);
    \draw[panelframe] (0,0) rectangle (10,8);
    \node[anchor=south west, font=\footnotesize, text=figink] at (0,8.2)
      {3\,$\cdot$\,lift: reroute seeking tails};
    \node[anchor=south west, font=\tiny, text=figgray] at (0.25,0.25)
      {under-served district (detail)};
    % recorded marks in the window (cell-snapped at 2x)
    \foreach \x/\y in {7/3, 5/3, 7/5, 9/3, 5/1, 7/7, 9/5, 9/7}
      \pic[color=figgray] at (\x,\y) {pickz};
    \pic at (5,5) {taxiz};
    % one trajectory, two futures. Shared approach (solid ink) to the anchor;
    % the original tail (solid gray) reaches its recorded pickup; the lift
    % reroute (dashed accent) bends into the highest-value cell.
    \draw[figink, line width=1.2pt] (0.6,7.6) -- (1.4,6.6) -- (2.2,5.8)
      -- (3,5);
    \draw[figgray, line width=1pt] (3,5) -- (3.8,3.6) -- (4.5,2.2) -- (5,1);
    \draw[liftdash] (3,5) -- (4.4,4.3) -- (5.8,3.6);
    \draw[liftdash, -{Stealth[length=4pt]}] (5.8,3.6) -- (7,3);
    \pic[color=figaccent] at (7,3) {pickz};
    % legend (upper right)
    \draw[figgray, line width=1pt] (6.1,7.5) -- (7.0,7.5);
    \node[anchor=west, font=\tiny, text=figink] at (7.15,7.5)
      {original tail};
    \draw[liftdash, -{Stealth[length=3.2pt]}] (6.1,6.9) -- (7.0,6.9);
    \node[anchor=west, font=\tiny, text=figink] at (7.15,6.9)
      {rerouted tail (lift)};
    \fill[figaccent, fill opacity=0.22] (6.1,6.1) rectangle (6.55,6.5);
    \node[anchor=west, font=\tiny, text=figink] at (7.15,6.3)
      {value of added presence};
  \end{scope}
\end{tikzpicture}
